\chapter{Méthode : une typologie textuelle du vocabulaire religieux}\label{ch:typologie}

\section{Pourquoi abandonner le profil unique}

Le catalogue de travail attribuait à chaque épitaphe un profil thématique unique (« aucune mention divine », « Dieu », « Christ », « Esprit saint », « binaire », « trinitaire explicite », « initiation et foi », « appartenance ecclésiale »), selon la catégorie la plus spécifique présente dans le texte. Ce codage avait le mérite de la clarté, mais il présentait deux défauts. Il ne distinguait pas le statut textuel des mots : un \emph{Christi} entièrement restitué par l’éditeur et un \emph{Chr(isti)} gravé en abrégé recevaient la même valeur. Et il projetait sur certains textes une structure doctrinale qu’ils ne contiennent pas : le \nnum{474}, qui dit seulement \emph{confirmans trinitatem}, était codé « trinitaire explicite (Père, Fils, Esprit) » ; le \nnum{458}, qui dit \emph{ancilla dei et Chr(isti)}, « binaire (Père et Fils) ».

La présente édition remplace ce profil par un relevé de chaque occurrence, qualifiée selon deux axes indépendants, et confronte ce relevé à une recherche lexicale indépendante menée sur l’ensemble du texte transmis.

\section{Premier axe : le statut textuel}

Quatre valeurs décrivent la transmission d’un mot : 1, mot lu ; 2, abréviation développée par l’éditeur ; 3, restitution partielle ; 4, restitution intégrale. Une marque distincte, S, désigne un signe figuré décrit par un témoin (monogramme du Christ). Deux codes hérités, 5 (« allusion ») et 6 (« anthroponyme »), décrivent une interprétation ou une fonction, non un état de conservation ; ils sont conservés pour la traçabilité, mais ne prolongent pas l’échelle philologique.

Une désignation divine n’est retenue indépendamment des restitutions que si l’élément conservé, ou une abréviation suffisamment déterminée, porte lui-même l’identification religieuse. Au \nnum{249}, \emph{spiritum} est lu, mais l’identification de l’Esprit saint dépend de la restitution partielle de \emph{[sanc]tum} ; au \nnum{352}, \emph{[qui Ch]ristum credens}, la fin du nom est lue mais son début restitué. Ces textes ne sont comptés qu’avec les restitutions.

Le niveau 1 ne certifie pas une lecture nouvellement collationnée sur la pierre : il signifie que les deux bases, ou la seule base disponible, transmettent le mot sans crochets. Lorsque l’EDB souligne des lettres, signe que ces lettres ont été lues par des éditeurs antérieurs mais ne sont plus visibles, la notice le signale (ainsi \emph{domino} au \nnum{85}).

\section{Second axe : l’emploi}

Chaque occurrence reçoit une catégorie d’emploi : désignation de Dieu, désignation du Christ, Esprit saint, Trinité, initiation et foi, vocabulaire ecclésial, saints, intercession, résurrection, repos et sommeil, allusion, anthroponyme. La catégorie dépend du contexte immédiat et a été attribuée par des règles explicites, contrôlées à la lecture : \emph{sanctum} suivi de \emph{spiritum} relève de l’Esprit saint ; \emph{sanctis} ou \emph{sanctorum} employés substantivement relèvent des saints ; \emph{sanctus} qualifiant un défunt est une épithète ; \emph{spiritus} désignant l’esprit du défunt relève de l’anthropologie. La titulature impériale \emph{dominus noster} est exclue, de même que \emph{dominus} au sens de maître (\emph{Cecili dominus Felicis}, \nnum{102}).

\section{Contrôle par une recherche indépendante}

Le relevé de la base de travail a été confronté à une recherche par expressions régulières sur le texte transmis de l’ensemble C, pour une trentaine de familles de mots (\emph{deus}, \emph{dominus}, \emph{Christus}, \emph{spiritus}, \emph{sanctus}, \emph{trinitas}, \emph{fides}, \emph{credere}, \emph{neophytus}, \emph{baptisma}, \emph{ecclesia}, \emph{martyr}, \emph{resurrectio}, \emph{caelum}, \emph{lux}, \emph{paradisus}, \emph{refrigerium}, \emph{anima}, \emph{gratia}, \emph{peccatum}, \emph{redemptus}, \emph{spes}, \emph{orare}…). Elle confirme le relevé pour les désignations divines. Elle fait apparaître des absences significatives : aucun texte latin de l’ensemble ne contient les mots \emph{Iesus}, \emph{baptisma}, \emph{martyr}, \emph{resurrectio}, \emph{paradisus}, \emph{gratia} ou \emph{peccatum}, même restitués. Elle conduit enfin à corriger trois classements :

\begin{itemize}
\item au \nnum{152}, la lecture \emph{redemptus}, qui fondait une allusion à la rédemption, n’est pas celle de l’EDB, qui lit \emph{[cor]p⟨o⟩re demptus} ; l’occurrence est retirée ;
\item au \nnum{447}, le mot \emph{ἀναστάσεως} est une restitution de l’EDCS, assortie d’un point d’interrogation, que l’EDB remplace par \emph{κρίσεως} ; il n’y a donc pas de mention de la résurrection assurée dans le corpus (chapitre~\ref{ch:mort_esperance}) ;
\item au \nnum{73}, la crainte de Dieu (\emph{[quaeque deu]m metuisti}) est une restitution de l’EDCS que l’EDB ne partage pas ; la désignation de Dieu du texte repose sur \emph{deo dica⟨vi⟩sti}, lu dans les deux bases.
\end{itemize}

\section{Résultats : trois profils de lecture}

Le tableau~\ref{tab:profils} réunit les dix-neuf textes de l’ensemble C qui contiennent, sous une forme quelconque, une désignation de Dieu, du Christ, de l’Esprit ou de la Trinité. Trois règles de compte en sont tirées.

\begin{description}[style=nextline,leftmargin=1.5em,font=\normalfont\scshape]
\item[Profil large] Textes où au moins une désignation est lue ou abrégée (niveaux 1–2), le signe du \nnum{38} exclu : \st{nProfilLarge} textes sur \st{nC} (\st{pProfilLarge}~\%).
\item[Profil restrictif] Le même, sans les deux textes dont l’identification dépend d’un développement douteux ou d’une restitution (\nnums{232} et 249) : \st{nProfilRestr} textes (\st{pProfilRestr}~\%).
\item[Profil strict] Le même, sans les trois textes où l’abréviation \emph{Chr.} rend probablement un monogramme (\nnums{140}, 151 et 311) : \st{nProfilStrict} textes (\st{pProfilStrict}~\%).
\end{description}

Avec les restitutions, le total atteint \st{nProfilAvecRest} textes (+ \nnums{43}, 104, 352, 382 et 447). Ces comptes décrivent des catégories de lecture, non un taux d’adhésion doctrinale. Ils montrent que le résultat principal — la rareté des désignations divines dans les épitaphes datées — ne dépend pas de la règle adoptée : entre 1,7 et 3,9~\% des textes selon la règle la plus sévère et la plus généreuse.

\begin{table}[tbp]\centering\small
\caption{Désignations de Dieu, du Christ, de l’Esprit et de la Trinité dans l’ensemble C. Statut : 1 lu, 2 abrégé, 3 restitution partielle, 4 restitution intégrale, S signe. Profils : L large, R restrictif, S strict ; « rest. » : compté seulement avec les restitutions.}\label{tab:profils}
\begin{tabularx}{\linewidth}{@{}r r X l l@{}}\toprule
N° & Année & Formes (statut) & Remarque & Profils\\\midrule
38 & 352 & \lat{in pace} + monogramme (S) & EDCS : \lat{Chr(isti)} & —\\
43 & 355 & \lat{in nomine [Christi]} (4) & chrisme décrit & rest.\\
73 & 363 & \lat{deo dica⟨vi⟩sti} (1) ; \lat{[deu]m} (3, EDCS seul) & restitutions divergentes & L R S\\
85 & 364 & \lat{omnia domino} (1) & EDB : lettres vues jadis & L R S\\
104 & 366 & \lat{in pace d[ei?]} (3) & EDB : \lat{d[- - -]} & rest.\\
140 & 370 & \lat{in pace Chr(isti)} (2) & chrisme décrit & L R\\
151 & 371 & \lat{in pace Chr(isti)} (2) & chrisme décrit & L R\\
214 & 380 & \lat{ancilla dei} (1) & & L R S\\
232 & 382 & \lat{in dom(ino?)} (2?) & développement douteux & L\\
249 & 383 & \lat{[sanc]tum spiritum} (1 + 3) & & L\\
306 & 389 & \lat{vitam Cristi}, \lat{ancilla dei} (1) ; \lat{[foedera C]hristi} (3) & ligne à contrôler & L R S\\
311 & 388 & \lat{in Chr(isto)} (2) & chrisme décrit & L R\\
352 & 393 & \lat{[qui Ch]ristum credens} (3) & & rest.\\
382 & 395 & \lat{in sin[u dei]} (4) & & rest.\\
400 & 397 & \lat{in luce domini} (1) & monogrammes décrits & L R S\\
447 & 401 & \lat{[θεὸν τὸν παντο]κράτορα} (4 + 3) & & rest.\\
458 & 401 & \lat{ancilla dei et Chr(ist)i} (1 + 2) & CHRI gravé & L R S\\
474 & 403 & \lat{homo dei}, \lat{trinitatem} (1) & & L R S\\
475 & 403 & \lat{dedit deus deus tulit} (1) & & L R S\\\bottomrule
\end{tabularx}
\end{table}

Ces dix-neuf textes sont examinés dans les chapitres suivants : Dieu et le Christ (chapitre~\ref{ch:dieu_christ}), l’Esprit et la Trinité (chapitre~\ref{ch:esprit_trinite}). Le vocabulaire de la foi, de l’initiation et de l’Église (chapitre~\ref{ch:foi_eglise}) et celui de la mort et de l’espérance (chapitre~\ref{ch:mort_esperance}) relèvent d’autres catégories d’emploi. Le répertoire du chapitre~\ref{ch:rep_religieux} donne toutes les occurrences, avec leur statut et leur contexte.
